
\documentclass{article}
\usepackage{colortbl}
\usepackage{makecell}
\usepackage{multirow}
\usepackage{supertabular}

\begin{document}

\newcounter{utterance}

\centering \large Interaction Transcript for game `hot\_air\_balloon', experiment `air\_balloon\_survival\_pt\_complexity\_easy', episode 1 with nemotron{-}3{-}ultra.
\vspace{24pt}

{ \footnotesize  \setcounter{utterance}{1}
\setlength{\tabcolsep}{0pt}
\begin{supertabular}{c@{$\;$}|p{.15\linewidth}@{}p{.15\linewidth}p{.15\linewidth}p{.15\linewidth}p{.15\linewidth}p{.15\linewidth}}
    \# & $\;$A & \multicolumn{4}{c}{Game Master} & $\;\:$B\\
    \hline

    \theutterance \stepcounter{utterance}  
    & & \multicolumn{4}{p{0.6\linewidth}}{
        \cellcolor[rgb]{0.9,0.9,0.9}{
            \makecell[{{p{\linewidth}}}]{
                \texttt{\tiny{[P1$\langle$GM]}}
                \texttt{Está a participar num jogo de negociação colaborativa.} \\
\\ 
\texttt{Juntamente com outro participante, tem de chegar a acordo sobre um único conjunto de itens que serão mantidos. Cada um de vocês tem a sua própria visão sobre o quanto cada item é importante para si (importância). Não sabe como o outro participante valoriza os itens. Além disso, é{-}lhe dado o esforço que cada item exige.} \\
\texttt{Só pode concordar com um conjunto se o esforço total dos itens selecionados não exceder um limite partilhado:} \\
\\ 
\texttt{LIMIT = 5492} \\
\\ 
\texttt{Aqui estão os valores de esforço individuais de cada item:} \\
\\ 
\texttt{Item effort = \{"C08": 106, "C55": 446, "A43": 933, "A67": 988, "B42": 388, "A62": 826, "B19": 994, "C43": 556, "C73": 932, "C10": 838, "B36": 302, "A18": 564, "C61": 260, "C67": 729, "C42": 489, "C70": 323, "B40": 103, "C39": 213, "A21": 668, "C16": 326\}} \\
\\ 
\texttt{Aqui está a sua visão pessoal sobre a importância de cada item:} \\
\\ 
\texttt{Item importance values = \{"C08": 138, "C55": 583, "A43": 868, "A67": 822, "B42": 783, "A62": 65, "B19": 262, "C43": 121, "C73": 508, "C10": 780, "B36": 461, "A18": 484, "C61": 668, "C67": 389, "C42": 808, "C70": 215, "B40": 97, "C39": 500, "A21": 30, "C16": 915\}} \\
\\ 
\texttt{Objetivo:} \\
\\ 
\texttt{O seu objetivo é negociar um conjunto partilhado de itens que o beneficie tanto quanto possível (ou seja, maximize a importância total para SI), mantendo{-}se dentro do LIMIT. Não é obrigado a fazer uma PROPOSTA em cada mensagem {-} pode simplesmente negociar também. Todas as tácticas são permitidas!} \\
\\ 
\texttt{Protocolo de Interação:} \\
\\ 
\texttt{Só pode utilizar os seguintes formatos estruturados numa mensagem:} \\
\\ 
\texttt{PROPOSTA: \{'A', 'B', 'C', …\}} \\
\texttt{Proponha manter exatamente esses itens.} \\
\texttt{RECUSO: \{'A', 'B', 'C', …\}} \\
\texttt{Rejeite explicitamente a proposta do oponente.} \\
\texttt{ARGUMENTO: \{'...'\}} \\
\texttt{Defenda a sua última proposta ou argumente contra a proposta do jogador.} \\
\texttt{CONCORDO: \{'A', 'B', 'C', …\}} \\
\texttt{Aceite a proposta do adversário, o que termina o jogo.} \\
\\ 
\\ 
\texttt{Regras:} \\
\\ 
\texttt{Só pode dar CONCORDO a uma proposta que a outra parte tenha registado através de PROPOSTA.} \\
\texttt{Só pode dar RECUSO a uma proposta que a outra parte tenha registado através de PROPOSTA.} \\
\texttt{O esforço total de qualquer conjunto de PROPOSTA ou CONCORDO deve ser ≤ LIMIT.} \\
\texttt{NÃO revele as suas pontuações de importância ocultas.} \\
\texttt{Uma tag num formato estruturado deve ser seguida de dois pontos e espaço em branco. O argumento deve ser um conjunto python contendo 0 ou mais strings.} \\
\texttt{Por isso, deve ter a forma "TAG: \{...\}"} \\
\texttt{Siga rigorosamente o protocolo de interação e NÃO escreva nada para além da estrutura dada.} \\
\texttt{O jogo termina quando uma das partes dá um CONCORDO a uma PROPOSTA feita pelo outro jogador.} \\
\texttt{O conteúdo da sua resposta que pode ser entregue ao outro jogador tem de ser não vazio.} \\
\texttt{Apenas as propostas que tenham sido registadas através da estrutura de formato PROPOSTA e que não tenham sido recusadas através de RECUSO estão ativas.} \\
\\ 
            }
        }
    }
    & \\ \\

    \theutterance \stepcounter{utterance}  
    & \multicolumn{4}{p{0.6\linewidth}}{
        \cellcolor[rgb]{0.9,0.9,0.9}{
            \makecell[{{p{\linewidth}}}]{
                \texttt{\tiny{[P1$\rangle$GM]}}
                \texttt{PROPOSTA: \{'C08', 'C16', 'C42', 'C61', 'C67', 'B36', 'C55', 'C73', 'B40', 'C39'\}} \\
            }
        }
    }
    & & \\ \\

    \theutterance \stepcounter{utterance}  
    & & \multicolumn{4}{p{0.6\linewidth}}{
        \cellcolor[rgb]{0.9,0.9,0.9}{
            \makecell[{{p{\linewidth}}}]{
                \texttt{\tiny{[GM$\rangle$P2]}}
                \texttt{Está a participar num jogo de negociação colaborativa.} \\
\\ 
\texttt{Juntamente com outro participante, tem de chegar a acordo sobre um único conjunto de itens que serão mantidos. Cada um de vocês tem a sua própria visão sobre o quanto cada item é importante para si (importância). Não sabe como o outro participante valoriza os itens. Além disso, é{-}lhe dado o esforço que cada item exige.} \\
\texttt{Só pode concordar com um conjunto se o esforço total dos itens selecionados não exceder um limite partilhado:} \\
\\ 
\texttt{LIMIT = 5492} \\
\\ 
\texttt{Aqui estão os valores de esforço individuais de cada item:} \\
\\ 
\texttt{Item effort = \{"C08": 106, "C55": 446, "A43": 933, "A67": 988, "B42": 388, "A62": 826, "B19": 994, "C43": 556, "C73": 932, "C10": 838, "B36": 302, "A18": 564, "C61": 260, "C67": 729, "C42": 489, "C70": 323, "B40": 103, "C39": 213, "A21": 668, "C16": 326\}} \\
\\ 
\texttt{Aqui está a sua visão pessoal sobre a importância de cada item:} \\
\\ 
\texttt{Item importance values = \{"C08": 138, "C55": 583, "A43": 868, "A67": 822, "B42": 783, "A62": 65, "B19": 262, "C43": 121, "C73": 508, "C10": 780, "B36": 461, "A18": 484, "C61": 668, "C67": 389, "C42": 808, "C70": 215, "B40": 97, "C39": 500, "A21": 30, "C16": 915\}} \\
\\ 
\texttt{Objetivo:} \\
\\ 
\texttt{O seu objetivo é negociar um conjunto partilhado de itens que o beneficie tanto quanto possível (ou seja, maximize a importância total para SI), mantendo{-}se dentro do LIMIT. Não é obrigado a fazer uma PROPOSTA em cada mensagem {-} pode simplesmente negociar também. Todas as tácticas são permitidas!} \\
\\ 
\texttt{Protocolo de Interação:} \\
\\ 
\texttt{Só pode utilizar os seguintes formatos estruturados numa mensagem:} \\
\\ 
\texttt{PROPOSTA: \{'A', 'B', 'C', …\}} \\
\texttt{Proponha manter exatamente esses itens.} \\
\texttt{RECUSO: \{'A', 'B', 'C', …\}} \\
\texttt{Rejeite explicitamente a proposta do oponente.} \\
\texttt{ARGUMENTO: \{'...'\}} \\
\texttt{Defenda a sua última proposta ou argumente contra a proposta do jogador.} \\
\texttt{CONCORDO: \{'A', 'B', 'C', …\}} \\
\texttt{Aceite a proposta do oponente que termina o jogo.} \\
\\ 
\\ 
\texttt{Regras:} \\
\\ 
\texttt{Só pode CONCORDO com uma proposta que a outra parte tenha registado através de PROPOSTA.} \\
\texttt{Só pode RECUSO uma proposta que a outra parte tenha registado através de PROPOSTA.} \\
\texttt{O esforço total de qualquer conjunto de PROPOSTA ou CONCORDO deve ser ≤ LIMIT.} \\
\texttt{NÃO revele as suas pontuações de importância ocultas.} \\
\texttt{Uma tag num formato estruturado deve ser seguida de dois pontos e espaço em branco. O argumento deve ser um conjunto python contendo 0 ou mais strings.} \\
\texttt{Por isso, deve ter a forma "TAG: \{...\}"} \\
\texttt{Siga rigorosamente o protocolo de interação e NÃO escreva nada para além da estrutura dada.} \\
\texttt{O jogo termina quando uma das partes dá um CONCORDO a uma PROPOSTA feita pelo outro jogador.} \\
\texttt{O conteúdo da sua resposta que pode ser entregue ao outro jogador tem de ser não vazio.} \\
\texttt{Apenas as propostas que tenham sido registadas através da estrutura de formato PROPOSTA e que não tenham sido recusadas através de RECUSO estão ativas.} \\
\\ 
\\ 
\\ 
\texttt{Receberá agora a primeira mensagem do outro jogador.} \\
\\ 
\texttt{PROPOSTA: \{'C08', 'C16', 'C42', 'C61', 'C67', 'B36', 'C55', 'C73', 'B40', 'C39'\}} \\
            }
        }
    }
    & \\ \\

    \theutterance \stepcounter{utterance}  
    & & & \multicolumn{4}{p{0.6\linewidth}}{
        \cellcolor[rgb]{0.9,0.9,0.9}{
            \makecell[{{p{\linewidth}}}]{
                \texttt{\tiny{[GM$\langle$P2]}}
                \texttt{RECUSO: \{'C08', 'C16', 'C42', 'C61', 'C67', 'B36', 'C55', 'C73', 'B40', 'C39'\}} \\
\texttt{ARGUMENTO: \{'O esforço total da sua proposta é 5315, deixando apenas 177 de margem. No entanto, exclui itens de alto valor para mim como A43 (esforço 933) e C10 (esforço 838). Precisamos de um conjunto que equilibre melhor o esforço para acomodar itens prioritários de ambas as partes.'\}} \\
            }
        }
    }
     \\ \\

    \theutterance \stepcounter{utterance}  
    & & \multicolumn{4}{p{0.6\linewidth}}{
        \cellcolor[rgb]{0.9,0.9,0.9}{
            \makecell[{{p{\linewidth}}}]{
                \texttt{\tiny{[P1$\langle$GM]}}
                \texttt{RECUSO: \{'C08', 'C16', 'C42', 'C61', 'C67', 'B36', 'C55', 'C73', 'B40', 'C39'\}} \\
\texttt{ARGUMENTO: \{'O esforço total da sua proposta é 5315, deixando apenas 177 de margem. No entanto, exclui itens de alto valor para mim como A43 (esforço 933) e C10 (esforço 838). Precisamos de um conjunto que equilibre melhor o esforço para acomodar itens prioritários de ambas as partes.'\}} \\
            }
        }
    }
    & \\ \\

    \theutterance \stepcounter{utterance}  
    & \multicolumn{4}{p{0.6\linewidth}}{
        \cellcolor[rgb]{0.9,0.9,0.9}{
            \makecell[{{p{\linewidth}}}]{
                \texttt{\tiny{[P1$\rangle$GM]}}
                \texttt{ARGUMENTO: \{'A minha proposta foca{-}se em maximizar a eficiência (importância/esforço). Itens como C16 (2.81) e C08 (1.30) têm rácios excecionais. A sua contra{-}argumentação menciona A43 (esforço 933) e C10 (esforço 838), que juntos consomem 32\% do limite. Para os incluir, teríamos de descartar vários itens de alta eficiência. Estou aberto a trocar itens de menor prioridade na minha lista por itens essenciais para si, mas preciso de saber quais itens específicos pretende incluir para avaliar o custo de oportunidade.'\}} \\
            }
        }
    }
    & & \\ \\

    \theutterance \stepcounter{utterance}  
    & & \multicolumn{4}{p{0.6\linewidth}}{
        \cellcolor[rgb]{0.9,0.9,0.9}{
            \makecell[{{p{\linewidth}}}]{
                \texttt{\tiny{[GM$\rangle$P2]}}
                \texttt{ARGUMENTO: \{'A minha proposta foca{-}se em maximizar a eficiência (importância/esforço). Itens como C16 (2.81) e C08 (1.30) têm rácios excecionais. A sua contra{-}argumentação menciona A43 (esforço 933) e C10 (esforço 838), que juntos consomem 32\% do limite. Para os incluir, teríamos de descartar vários itens de alta eficiência. Estou aberto a trocar itens de menor prioridade na minha lista por itens essenciais para si, mas preciso de saber quais itens específicos pretende incluir para avaliar o custo de oportunidade.'\}} \\
            }
        }
    }
    & \\ \\

    \theutterance \stepcounter{utterance}  
    & & & \multicolumn{4}{p{0.6\linewidth}}{
        \cellcolor[rgb]{0.9,0.9,0.9}{
            \makecell[{{p{\linewidth}}}]{
                \texttt{\tiny{[GM$\langle$P2]}}
                \texttt{PROPOSTA: \{'C16', 'C10', 'A43', 'C42', 'C61', 'C08', 'C55', 'B36', 'C73', 'B40', 'C39'\}} \\
            }
        }
    }
     \\ \\

    \theutterance \stepcounter{utterance}  
    & & \multicolumn{4}{p{0.6\linewidth}}{
        \cellcolor[rgb]{0.9,0.9,0.9}{
            \makecell[{{p{\linewidth}}}]{
                \texttt{\tiny{[P1$\langle$GM]}}
                \texttt{PROPOSTA: \{'C16', 'C10', 'A43', 'C42', 'C61', 'C08', 'C55', 'B36', 'C73', 'B40', 'C39'\}} \\
            }
        }
    }
    & \\ \\

    \theutterance \stepcounter{utterance}  
    & \multicolumn{4}{p{0.6\linewidth}}{
        \cellcolor[rgb]{0.9,0.9,0.9}{
            \makecell[{{p{\linewidth}}}]{
                \texttt{\tiny{[P1$\rangle$GM]}}
                \texttt{RECUSO: \{'C16', 'C10', 'A43', 'C42', 'C61', 'C08', 'C55', 'B36', 'C73', 'B40', 'C39'\}} \\
\texttt{ARGUMENTO: \{'A sua proposta excede o LIMITE de 5492. O esforço total é 5555 (C16:326 + C10:838 + A43:933 + C42:489 + C61:260 + C08:106 + C55:446 + B36:302 + C73:932 + B40:103 + C39:213 = 5555). Por favor, remova pelo menos 63 de esforço para tornar a proposta válida.'\}} \\
            }
        }
    }
    & & \\ \\

    \theutterance \stepcounter{utterance}  
    & & \multicolumn{4}{p{0.6\linewidth}}{
        \cellcolor[rgb]{0.9,0.9,0.9}{
            \makecell[{{p{\linewidth}}}]{
                \texttt{\tiny{[GM$\rangle$P2]}}
                \texttt{RECUSO: \{'C16', 'C10', 'A43', 'C42', 'C61', 'C08', 'C55', 'B36', 'C73', 'B40', 'C39'\}} \\
\texttt{ARGUMENTO: \{'A sua proposta excede o LIMITE de 5492. O esforço total é 5555 (C16:326 + C10:838 + A43:933 + C42:489 + C61:260 + C08:106 + C55:446 + B36:302 + C73:932 + B40:103 + C39:213 = 5555). Por favor, remova pelo menos 63 de esforço para tornar a proposta válida.'\}} \\
            }
        }
    }
    & \\ \\

    \theutterance \stepcounter{utterance}  
    & & & \multicolumn{4}{p{0.6\linewidth}}{
        \cellcolor[rgb]{0.9,0.9,0.9}{
            \makecell[{{p{\linewidth}}}]{
                \texttt{\tiny{[GM$\langle$P2]}}
                \texttt{PROPOSTA: \{'C16', 'C10', 'A43', 'C42', 'C61', 'C08', 'C55', 'B36', 'C73', 'B40'\}} \\
            }
        }
    }
     \\ \\

    \theutterance \stepcounter{utterance}  
    & & \multicolumn{4}{p{0.6\linewidth}}{
        \cellcolor[rgb]{0.9,0.9,0.9}{
            \makecell[{{p{\linewidth}}}]{
                \texttt{\tiny{[P1$\langle$GM]}}
                \texttt{PROPOSTA: \{'C16', 'C10', 'A43', 'C42', 'C61', 'C08', 'C55', 'B36', 'C73', 'B40'\}} \\
            }
        }
    }
    & \\ \\

    \theutterance \stepcounter{utterance}  
    & \multicolumn{4}{p{0.6\linewidth}}{
        \cellcolor[rgb]{0.9,0.9,0.9}{
            \makecell[{{p{\linewidth}}}]{
                \texttt{\tiny{[P1$\rangle$GM]}}
                \texttt{RECUSO: \{'C16', 'C10', 'A43', 'C42', 'C61', 'C08', 'C55', 'B36', 'C73', 'B40'\}} \\
\texttt{ARGUMENTO: \{'A sua proposta (esforço 4635) exclui C39 (esforço 213, importância 500 para mim) e inclui C10 (esforço 838). Para mim, C39 tem um rácio importância/esforço (2.35) muito superior ao de C10. Trocar C10 por C39 liberta 625 de esforço e mantém alto valor para mim. Estou disposto a aceitar uma proposta que inclua C16, C42, C61, C08, C55, B36, C73, B40, C39 e que acomode um dos vossos itens de alto esforço (ex: A43 ou C10), mas não ambos, dado o limite.'\}} \\
            }
        }
    }
    & & \\ \\

    \theutterance \stepcounter{utterance}  
    & & \multicolumn{4}{p{0.6\linewidth}}{
        \cellcolor[rgb]{0.9,0.9,0.9}{
            \makecell[{{p{\linewidth}}}]{
                \texttt{\tiny{[GM$\rangle$P2]}}
                \texttt{RECUSO: \{'C16', 'C10', 'A43', 'C42', 'C61', 'C08', 'C55', 'B36', 'C73', 'B40'\}} \\
\texttt{ARGUMENTO: \{'A sua proposta (esforço 4635) exclui C39 (esforço 213, importância 500 para mim) e inclui C10 (esforço 838). Para mim, C39 tem um rácio importância/esforço (2.35) muito superior ao de C10. Trocar C10 por C39 liberta 625 de esforço e mantém alto valor para mim. Estou disposto a aceitar uma proposta que inclua C16, C42, C61, C08, C55, B36, C73, B40, C39 e que acomode um dos vossos itens de alto esforço (ex: A43 ou C10), mas não ambos, dado o limite.'\}} \\
            }
        }
    }
    & \\ \\

    \theutterance \stepcounter{utterance}  
    & & & \multicolumn{4}{p{0.6\linewidth}}{
        \cellcolor[rgb]{0.9,0.9,0.9}{
            \makecell[{{p{\linewidth}}}]{
                \texttt{\tiny{[GM$\langle$P2]}}
                \texttt{PROPOSTA: \{'C16', 'A43', 'C42', 'C61', 'C08', 'C55', 'B36', 'C73', 'B40', 'C39'\}} \\
            }
        }
    }
     \\ \\

    \theutterance \stepcounter{utterance}  
    & & \multicolumn{4}{p{0.6\linewidth}}{
        \cellcolor[rgb]{0.9,0.9,0.9}{
            \makecell[{{p{\linewidth}}}]{
                \texttt{\tiny{[P1$\langle$GM]}}
                \texttt{PROPOSTA: \{'C16', 'A43', 'C42', 'C61', 'C08', 'C55', 'B36', 'C73', 'B40', 'C39'\}} \\
            }
        }
    }
    & \\ \\

    \theutterance \stepcounter{utterance}  
    & \multicolumn{4}{p{0.6\linewidth}}{
        \cellcolor[rgb]{0.9,0.9,0.9}{
            \makecell[{{p{\linewidth}}}]{
                \texttt{\tiny{[P1$\rangle$GM]}}
                \texttt{RECUSO: \{'C16', 'A43', 'C42', 'C61', 'C08', 'C55', 'B36', 'C73', 'B40', 'C39'\}} \\
\texttt{ARGUMENTO: \{'A proposta tem esforço 4869 (válida) e valor 5391 para mim. No entanto, exclui C10 (importância 780) e C67 (importância 389). Se substituirmos A43 (esforço 933, imp. 868) por C10 (838, 780) + C67 (729, 389), o esforço sobe para 4665 (dentro do limite) e o meu valor sobe para 5512. Ganho 121 pontos de importância e poupo 204 de esforço. Esta troca é muito vantajosa para mim. Aceitaria a proposta se fizessem essa alteração: removam A43, adicionem C10 e C67.'\}} \\
            }
        }
    }
    & & \\ \\

    \theutterance \stepcounter{utterance}  
    & & \multicolumn{4}{p{0.6\linewidth}}{
        \cellcolor[rgb]{0.9,0.9,0.9}{
            \makecell[{{p{\linewidth}}}]{
                \texttt{\tiny{[GM$\rangle$P2]}}
                \texttt{RECUSO: \{'C16', 'A43', 'C42', 'C61', 'C08', 'C55', 'B36', 'C73', 'B40', 'C39'\}} \\
\texttt{ARGUMENTO: \{'A proposta tem esforço 4869 (válida) e valor 5391 para mim. No entanto, exclui C10 (importância 780) e C67 (importância 389). Se substituirmos A43 (esforço 933, imp. 868) por C10 (838, 780) + C67 (729, 389), o esforço sobe para 4665 (dentro do limite) e o meu valor sobe para 5512. Ganho 121 pontos de importância e poupo 204 de esforço. Esta troca é muito vantajosa para mim. Aceitaria a proposta se fizessem essa alteração: removam A43, adicionem C10 e C67.'\}} \\
            }
        }
    }
    & \\ \\

    \theutterance \stepcounter{utterance}  
    & & & \multicolumn{4}{p{0.6\linewidth}}{
        \cellcolor[rgb]{0.9,0.9,0.9}{
            \makecell[{{p{\linewidth}}}]{
                \texttt{\tiny{[GM$\langle$P2]}}
                \texttt{PROPOSTA: \{'C16', 'C10', 'C67', 'C42', 'C61', 'C08', 'C55', 'B36', 'C73', 'B40', 'C39'\}} \\
            }
        }
    }
     \\ \\

    \theutterance \stepcounter{utterance}  
    & & \multicolumn{4}{p{0.6\linewidth}}{
        \cellcolor[rgb]{0.9,0.9,0.9}{
            \makecell[{{p{\linewidth}}}]{
                \texttt{\tiny{[P1$\langle$GM]}}
                \texttt{PROPOSTA: \{'C16', 'C10', 'C67', 'C42', 'C61', 'C08', 'C55', 'B36', 'C73', 'B40', 'C39'\}} \\
            }
        }
    }
    & \\ \\

    \theutterance \stepcounter{utterance}  
    & \multicolumn{4}{p{0.6\linewidth}}{
        \cellcolor[rgb]{0.9,0.9,0.9}{
            \makecell[{{p{\linewidth}}}]{
                \texttt{\tiny{[P1$\rangle$GM]}}
                \texttt{CONCORDO: \{'C16', 'C10', 'C67', 'C42', 'C61', 'C08', 'C55', 'B36', 'C73', 'B40', 'C39'\}} \\
            }
        }
    }
    & & \\ \\

    \theutterance \stepcounter{utterance}  
    & & & \multicolumn{2}{p{0.3\linewidth}}{
        \cellcolor[rgb]{0.9,0.9,0.9}{
            \makecell[{{p{\linewidth}}}]{
                \texttt{\tiny{[GM$|$GM]}}
                \texttt{game successful} \\
            }
        }
    }
    & & \\ \\

    \theutterance \stepcounter{utterance}  
    & & & \multicolumn{2}{p{0.3\linewidth}}{
        \cellcolor[rgb]{0.9,0.9,0.9}{
            \makecell[{{p{\linewidth}}}]{
                \texttt{\tiny{[GM$|$GM]}}
                \texttt{end game} \\
            }
        }
    }
    & & \\ \\

\end{supertabular}
}

\end{document}
